\documentclass{article}
\usepackage[UTF8]{ctex}
\usepackage{amsmath,amssymb,amsfonts}
\usepackage{enumitem}
\usepackage{geometry}
\geometry{a4paper, margin=1in}
% === 错题本排版优化（不参与题目标识） ===
\setlength{\parindent}{0pt}
\setlength{\parskip}{0.35em plus 0.10em minus 0.08em}
\setlength{\emergencystretch}{3em}
\setlist{itemsep=0.25em, topsep=0.40em, parsep=0pt}
\allowdisplaybreaks[2]
% === 错题本排版优化结束 ===
\begin{document}
\section*{北京科技大学2026年数学分析真题}
\section*{13．北京科技大学2026年数学分析试卷}

\subsection{题目1}
1．求极限 \(\lim _{x \rightarrow\left(\frac{\pi}{2}\right)^{-}}(\cos x)^{\frac{\pi}{2}-x}\) ．

\subsection{题目2}
2．设 \(a=\sum_{i=1}^{n} a_{i}, a_{i}>0\) ，求 \(S=\sum_{i=1}^{n} a_{i}^{2}\) 的最小值．

\subsection{题目3}
3．求积分 \(\int_{\frac{1}{2}}^{\frac{3}{4}} \frac{\arcsin \sqrt{x}}{\sqrt{x(1-x)}} \mathrm{d} x\) ．

\subsection{题目4}
4．设 \(\left\{x_{n}\right\}\) 满足 \(x_{1}=1, x_{n+1}=x_{n}+\frac{1}{2 x_{n}}(n=1,2, \cdots)\) ，求极限 \(\lim _{n \rightarrow \infty} \frac{x_{n}^{2}-n}{\ln n}\) ．

二．设 \(f(x)\) 在 \([0,+\infty)\) 上连续可微， \(\lim _{x \rightarrow+\infty} f^{\prime}(x)=A(A\) 为常数 \()\) ，证明： \(\lim _{x \rightarrow+\infty} \frac{f(x)}{x}=A\) ．

三．已知 \(f(x, y)=\varphi(|x y|)\) ，其中 \(\varphi(0)=0\) ，且 \(u=0\) 附近满足 \(|\varphi(u)| \leq u^{2}\) ，证明：\(f(x, y)\) 在 \((0,0)\)处可微．

四．利用有限覆盖定理证明连续函数的根的存在定理．

五．求级数 \(\sum_{n=1}^{\infty} \frac{(-1)^{n} a^{n}}{n \ln \left(n^{2}+n\right)} x^{2 n-1}\) 的收敛域，其中 \(a>1\) 。

六．对于级数 \(\sum_{n=1}^{\infty} \frac{\ln \left(1+n^{2} x\right)}{n} x^{-(n+1)}\) ．证明：

\begin{itemize}
  \item[（1）] 该级数在 \(\left[1+\varepsilon_{0},+\infty\right)\left(\varepsilon_{0}>0\right)\) 上一致收敛．
  \item[（2）] 该级数在 \((1,+\infty)\) 上不一致收敛．
\end{itemize}

\subsection{题目7}
原题边界“七．”：讨论 \(I(a)=\int_{0}^{+\infty} \frac{a}{1+x^{2} a^{2}} \mathrm{~d} x\) 在 \(\left[a_{0},+\infty\right)\) 和 \((0,+\infty)\) 的一致收敛性，其中 \(a_{0}>0\) ．

八．设 \(f(x)\) 是可导函数，对任意围绕原点一周的定向闭曲线 \(L\) ，曲线积分 \(\int_{L} \frac{x \mathrm{~d} y-y \mathrm{~d} x}{f(x)+4 y^{2}}\) 的值恒为常数\(A\) ，且 \(f(1)=2\) ，求 \(A\) 的值与 \(f(x)\) ．

九．讨论 \(\int_{0}^{+\infty} \frac{x^{p-1}}{1+x} \mathrm{~d} x\) 的敛散性．
\end{document}
